\chapter{Fundamentos Analíticos y Formalismo de la Zeta-Regularización}

\section{El problema de la divergencia}

El análisis de las series infinitas y los productos infinitos ha sido, desde los tiempos de Euler, una fuente inagotable de paradojas y descubrimientos profundos. En el contexto de la Zeta-Regularización, el \textit{problema de la divergencia} no es meramente un inconveniente técnico que debe ser eludido, sino la puerta de entrada a una estructura matemática más rica que conecta el análisis complejo, la teoría de números y la física teórica. Para comprender la necesidad y la potencia de la regularización, primero debemos diseccionar rigurosamente por qué las sumas y productos clásicos fracasan al evaluar secuencias no acotadas.

\subsection{Series de Dirichlet y la abscisa de convergencia}

El punto de partida natural para estudiar sumas infinitas de términos que crecen o decrecen polinómicamente es la Serie de Dirichlet. En el contexto de la zeta-regularización, la secuencia $\Lambda = \{\lambda_n\}_{n=1}^{\infty}$ representa típicamente el espectro de un operador diferencial elíptico positivo (como el Laplaciano en una variedad de Riemann) o un conjunto aritmético específico (como los números naturales, los números primos o los números odiosos).

A cada secuencia admisible se le asocia una función zeta espectral $\zeta_\Lambda(s)$, cuya definición rigurosa se presentará en la Sección siguiente. Por ahora, basta observar que su comportamiento está gobernado por una \textit{abscisa de convergencia} $\sigma_c \in \mathbb{R} \cup \{-\infty, \infty\}$: la serie converge absolutamente para $\Re(s) > \sigma_c$ y diverge para $\Re(s) < \sigma_c$.

La abscisa $\sigma_c$ está intrínsecamente ligada a la tasa de crecimiento $\alpha$ de la función de conteo $N(x)$. Por ejemplo:
\begin{itemize}
    \item Si $\Lambda = \mathbb{N} = \{1, 2, 3, \dots\}$, la secuencia representa los números naturales. La función resultante es la función Zeta de Riemann, $\zeta(s) = \sum n^{-s}$, la cual tiene una abscisa de convergencia $\sigma_c = 1$.
    \item Si $\Lambda = \mathbb{P} = \{2, 3, 5, 7, \dots\}$, la secuencia representa los números primos. La función resultante es la función Zeta Prima, $P(s) = \sum p^{-s}$, que también posee $\sigma_c = 1$.
    \item Si $\Lambda = \{\text{valores propios de } \Delta\}$ en una variedad compacta de dimensión $d$, la ley de Weyl asegura que $N(x) \sim C x^{d/2}$, lo que implica que la abscisa de convergencia es $\sigma_c = d/2$. En este contexto físico-geométrico, $\Delta$ es el operador de Laplace-Beltrami sobre la variedad de Riemann, y sus valores propios $\lambda_n$ representan los modos de vibración o niveles de energía del sistema. La ley de Weyl establece que el número de valores propios menores o iguales a $x$ crece asintóticamente como $C x^{d/2}$, donde $C$ depende del volumen de la variedad. Esto significa que los valores propios crecen como $\lambda_n \sim n^{2/d}$, y por tanto la serie de Dirichlet $\sum \lambda_n^{-s}$ converge absolutamente cuando $\Re(s) > d/2$. Por ejemplo, para una superficie bidimensional ($d=2$), la abscisa de convergencia es $\sigma_c = 1$, mientras que para un dominio tridimensional ($d=3$), es $\sigma_c = 3/2$. Esta estructura es fundamental en geometría espectral y en la regularización de determinantes funcionales en teoría cuántica de campos.
\end{itemize}

El problema fundamental de la divergencia surge cuando deseamos evaluar la serie, o sus derivadas, en puntos donde $\Re(s) \le \sigma_c$. Por ejemplo, si $\Lambda = \mathbb{N}$ e intentamos evaluar la suma de todos los números naturales, nos enfrentamos a la serie $\sum_{n=1}^{\infty} n$, que corresponde a evaluar $\zeta_\Lambda(s)$ en $s = -1$. Clásicamente, esta serie diverge hacia $+\infty$. De manera similar, si $\Lambda$ representa los valores propios de un operador y queremos calcular el producto de todos ellos (el determinante funcional), el valor regularizado relevante queda ligado al comportamiento de $\zeta_\Lambda(s)$ cerca de $s = 0$. Si la abscisa de convergencia es $\sigma_c > 0$ (como ocurre en casi todos los casos físicos y aritméticos de interés), el punto $s=0$ se encuentra en la región de divergencia clásica, haciendo que la evaluación directa de la serie sea matemáticamente imposible bajo los criterios estándar del análisis real.

\subsection{La ilusión del producto infinito y la suma logarítmica}

El problema se intensifica drásticamente cuando pasamos de las sumas a los productos infinitos. Dada una secuencia $\{\lambda_n\}_{n=1}^{\infty}$, el producto infinito se define formalmente como:

$$ P = \prod_{n=1}^{\infty} \lambda_n $$

En el análisis clásico, para que un producto infinito $\prod (1 + a_n)$ converja a un valor distinto de cero, es una condición necesaria que $\lim_{n \to \infty} a_n = 0$. En nuestro caso, si $\lambda_n \to \infty$, el producto diverge trivialmente hacia el infinito. La intuición clásica nos dice que multiplicar infinitos números mayores que 1 (eventualmente) debe producir una magnitud infinita.

Sin embargo, la conexión formal entre productos y sumas a través del logaritmo:

$$ \log \left( \prod_{n=1}^{N} \lambda_n \right) = \sum_{n=1}^{N} \log \lambda_n $$

nos revela la naturaleza de la divergencia. Si tomamos el límite $N \to \infty$, la suma $\sum \log \lambda_n$ diverge. Por ejemplo, para el producto de todos los números naturales $P = 1 \cdot 2 \cdot 3 \cdots$, la suma logarítmica es $\sum_{n=1}^{\infty} \log n$, la cual diverge logarítmicamente. 

No obstante, Euler y otros matemáticos notaron que ciertas identidades formales, como la manipulación de series de potencias o la fórmula de los productos de Weierstrass, parecían requerir que estos productos infinitos tuvieran valores finitos y bien definidos. Euler asignó el valor $\sqrt{2\pi}$ al producto de todos los números naturales, un resultado que carecía de justificación en el análisis clásico del siglo XVIII, pero que anticipaba la necesidad de un nuevo paradigma.

\subsection{Motivación en Física Teórica y Geometría Espectral}

La necesidad de resolver el problema de la divergencia no es solo una curiosidad analítica; es una exigencia física y geométrica. En múltiples áreas de la ciencia, las cantidades observables están definidas por productos o sumas sobre todos los modos de un sistema, los cuales son infinitos y divergentes.

\begin{itemize}
    \item \textbf{Energía del punto cero y el Efecto Casimir:} En la Teoría Cuántica de Campos, el vacío no está "vacío", sino que fluctúa. La energía del vacío para un campo escalar en una región del espacio se calcula sumando las energías de punto cero de todos los modos electromagnéticos permitidos: $E_0 = \frac{1}{2} \sum_n \hbar \omega_n$. Para el efecto Casimir entre dos placas paralelas separadas por una distancia $d$, las frecuencias están cuantizadas y la suma resultante es proporcional a $\sum_{n=1}^{\infty} n^3$ o $\sum_{n=1}^{\infty} n$, dependiendo de la dimensionalidad y la geometría. Clásicamente, esta energía es infinita, pero físicamente solo importa la \textit{diferencia} de energía entre la configuración con placas y el vacío sin restricciones.
    \item \textbf{Determinantes Funcionales en Geometría Espectral:} En geometría diferencial y teoría de cuerdas, a menudo se necesita calcular el determinante de un operador diferencial elíptico, como el operador de Laplace-Beltrami $\Delta_g$ en una variedad de Riemann $(M, g)$. Los valores propios $\{\lambda_n\}$ del Laplace-Beltrami crecen asintóticamente (ley de Weyl), por lo que el producto $\det(\Delta_g) = \prod_{n=1}^{\infty} \lambda_n$ diverge catastróficamente. Este determinante es crucial para definir la medida de integración en integrales de camino (path integrals) de Hawking y para calcular la torsión de Reidemeister analítica (torsión de Ray-Singer).
    \item \textbf{Magnetización en Materiales Cuánticos:} En sistemas de materia condensada como el grafeno bajo un campo magnético fuerte, el potencial termodinámico requiere sumar sobre los infinitos Niveles de Landau. La suma $\sum \sqrt{n}$ diverge, pero la respuesta física observable (la magnetización) es finita y medible experimentalmente.
\end{itemize}

En todos estos casos, la física nos grita que la cantidad subyacente debe ser finita. El problema de la divergencia es, por tanto, un artefacto de aplicar herramientas de convergencia clásica a problemas que poseen una estructura analítica subyacente más profunda.

\subsection{El enfoque de la regularización: más allá de la convergencia clásica}

Para resolver el problema de la divergencia, las matemáticas modernas introducen el concepto de \textit{regularización}. Una regularización no es un "truco" para ignorar los infinitos, sino un procedimiento sistemático y riguroso para extender el dominio de definición de un funcional. 

Formalmente, buscamos un método que asigne un valor finito a una suma o producto divergente satisfaciendo ciertas propiedades de consistencia. La Zeta-Regularización logra esto sin eliminar los términos infinitos: cambia el problema de una evaluación directa de series divergentes a una evaluación analítica de una función compleja continuada fuera de su dominio inicial de convergencia.

El problema de la divergencia, por lo tanto, se resuelve no eliminando los términos infinitos, sino redefiniendo el espacio de funciones en el que operamos, utilizando el plano complejo como un lienzo donde las singularidades pueden ser navegadas mediante la continuación analítica, revelando el valor finito y físico que yace oculto debajo de la divergencia clásica.


\section{Definición formal del producto y suma Zeta-regularizada}

El corazón de la Zeta-Regularización reside en transformar operaciones aritméticas clásicamente divergentes (sumas y productos infinitos) en operaciones analíticas bien definidas sobre el plano complejo. Para ello, no evaluamos la suma o el producto directamente, sino que construimos un objeto analítico intermediario —la función zeta espectral— y extraemos los valores de interés mediante técnicas de continuación analítica y derivación compleja.

\subsection{Construcción de la función zeta asociada a una secuencia}

Sea $\Lambda = \{\lambda_n\}_{n=1}^{\infty}$ una secuencia de números complejos no nulos. Para que la construcción analítica sea viable, imponemos tres condiciones de crecimiento y ordenamiento:

\begin{enumerate}[a)]
    \item \textbf{Ordenamiento por magnitud y ausencia de acumulación:} Los elementos están ordenados de modo que $0 < |\lambda_1| \le |\lambda_2| \le \dots$. El cero queda excluido ($0 \notin \Lambda$), pues los términos $\lambda_n^{-s}$ serían indefinidos si algún $\lambda_n$ fuera nulo. Además, $\Lambda$ es discreto y carece de puntos de acumulación finitos: para todo $R>0$, el disco $\{z \in \mathbb{C}: |z| \le R\}$ contiene solo un número finito de elementos de $\Lambda$.
    \item \textbf{Divergencia hacia el infinito:} La magnitud de los elementos crece indefinidamente, es decir, $\lim_{n \to \infty} |\lambda_n| = \infty$.
    \item \textbf{Cota de crecimiento polinomial:} La función de conteo espectral
    $$ N(x)=|\{n\in\mathbb{N}: |\lambda_n|\le x\}| $$
    satisface $N(x)=\mathcal{O}(x^\alpha)$ para alguna constante $\alpha>0$. Esta hipótesis garantiza que la serie de Dirichlet asociada tenga al menos una región de convergencia absoluta en un semiplano derecho.
\end{enumerate}

Para construir la función zeta asociada a esta secuencia, definimos la Serie de Dirichlet generalizada:

$$ \zeta_\Lambda(s) = \sum_{n=1}^{\infty} \frac{1}{\lambda_n^s} = \sum_{n=1}^{\infty} \lambda_n^{-s} = \sum_{\lambda \in \Lambda} \lambda^{-s} $$

donde $s = \sigma + it \in \mathbb{C}$ es una variable compleja.

La construcción de $\zeta_\Lambda(s)$ no es arbitraria; responde a una necesidad estructural. Al elevar los elementos de la secuencia a la potencia $-s$, convertimos el crecimiento asintótico de $\lambda_n$ en un mecanismo de decaimiento exponencial para la parte real de $s$ suficientemente grande. Específicamente, existe una abscisa de convergencia $\sigma_c$ tal que la serie converge absolutamente y define una función holomorfa en el semiplano derecho $\Re(s) > \sigma_c$.

En este dominio de convergencia, la función $\zeta_\Lambda(s)$ codifica toda la información asintótica y aritmética de la secuencia $\Lambda$. Sin embargo, como discutimos previamente, los valores de interés físico y matemático (como la suma de los elementos o el producto de los mismos) corresponden a evaluaciones en $s = -1$ o $s = 0$, puntos que yacen en la región de divergencia clásica $\Re(s) \le \sigma_c$.

\subsection{Continuación analítica y evaluación en $s=0$}

Dado que la serie de Dirichlet original diverge en $s=0$, debemos invocar el principio de continuación analítica. Asumimos que $\zeta_\Lambda(s)$ admite una continuación meromorfa a todo el plano complejo $\mathbb{C}$, o al menos a un entorno abierto que contenga el origen $s=0$. 

El comportamiento de la función en $s=0$ dicta la definición formal del producto regularizado. Aquí nos enfrentamos a dos escenarios analíticos distintos:

\begin{enumerate}[a)]
    \item \textbf{El caso regular (Origen sin polo):} Si la continuación meromorfa de $\zeta_\Lambda(s)$ es holomorfa en $s=0$ (es decir, $s=0$ es un punto regular), entonces la función es analítica en un entorno del origen. En este caso, podemos evaluar directamente la función y sus derivadas en $s=0$. El valor $\zeta_\Lambda(0)$ estará relacionado con la regularización de la suma de los logaritmos, y la derivada $\zeta_\Lambda'(0)$ será la clave para el producto regularizado.
    
    \item \textbf{El caso singular (Origen con polo):} En muchas aplicaciones avanzadas (como en ciertos productos de secuencias automáticas o funciones zeta trigonométricas sobre ceros de Riemann), la función $\zeta_\Lambda(s)$ puede presentar un polo en $s=0$. Si $s=0$ es un polo, la derivada estándar $\zeta_\Lambda'(0)$ no está definida. Para resolver esto, se introduce una generalización de la regularización. Sea $P(s; \Lambda)$ la parte meromorfa de $\zeta_\Lambda(s)$ en $s=0$. Definimos el \textit{término lineal} de $\zeta_\Lambda(s)$ en el origen como el residuo:
    $$ \text{LT}_{s=0} \zeta_\Lambda(s) = \text{Res}_{s=0} \frac{\zeta_\Lambda(s)}{s^2} $$
    Este residuo extrae el coeficiente del término lineal en la expansión de Laurent de $\zeta_\Lambda(s)$ alrededor de $s=0$, sustituyendo eficazmente el papel de $\zeta_\Lambda'(0)$ cuando hay un polo.
\end{enumerate}

Esta distinción es crucial. Mientras que en la física espectral estándar (como el Laplaciano en variedades compactas) $s=0$ suele ser un punto regular, en teoría de números analítica y geometría hiperbólica es frecuente encontrarse con polos en el origen, requiriendo el uso de esta definición generalizada basada en residuos.

\subsection{La fórmula maestra: Sumas y Productos Regularizados}

Una vez que la función $\zeta_\Lambda(s)$ ha sido continuada analíticamente y su comportamiento en $s=0$ ha sido clasificado, podemos definir rigurosamente las sumas y productos regularizados. 

\subsubsection*{Derivación formal de la fórmula maestra}

Para comprender la intuición detrás de la fórmula, consideremos la manipulación formal (heurística) del logaritmo de un producto infinito. Si asumimos que las propiedades de los logaritmos se mantienen para productos infinitos, tenemos:

$$ \log \left( \prod_{n=1}^{\infty} \lambda_n \right) = \sum_{n=1}^{\infty} \log \lambda_n $$

Observamos que la derivada de la función zeta espectral respecto a $s$ es:

$$ \zeta_\Lambda'(s) = \frac{d}{ds} \sum_{n=1}^{\infty} \lambda_n^{-s} = \frac{d}{ds} \sum_{n=1}^{\infty} e^{-s \log \lambda_n} = -\sum_{n=1}^{\infty} (\log \lambda_n) \lambda_n^{-s} $$

Si evaluamos esta derivada formalmente en $s=0$, obtenemos:

$$ \zeta_\Lambda'(0) = -\sum_{n=1}^{\infty} \log \lambda_n = -\log \left( \prod_{n=1}^{\infty} \lambda_n \right) $$

Despejando el producto, llegamos a la expresión formal $\prod \lambda_n = \exp(-\zeta_\Lambda'(0))$. Sin embargo, dado que la serie para $\zeta_\Lambda'(0)$ diverge en el análisis clásico, la Zeta-Regularización eleva esta heurística a una definición rigurosa mediante la continuación analítica.

\subsubsection*{Definición rigurosa}

Basándonos en el análisis anterior, definimos los siguientes operados regularizados:

\begin{itemize}
    \item \textbf{Producto Zeta-regularizado (Caso regular):} Si $\zeta_\Lambda(s)$ es holomorfa en $s=0$, el producto regularizado de la secuencia $\Lambda$ se define mediante la \textbf{fórmula maestra}:
    $$ \prod_{n=1}^{\infty} \lambda_n := \exp(-\zeta_\Lambda'(0)) $$
    Esta es la definición estándar utilizada en geometría espectral para calcular el determinante funcional de un operador, $\det(\Delta) = \exp(-\zeta_\Delta'(0))$.

    \item \textbf{Producto Zeta-regularizado (Caso generalizado con polo):} Si $\zeta_\Lambda(s)$ tiene un polo en $s=0$, el producto regularizado (a menudo llamado \textit{producto ddotted} en la literatura especializada) se define utilizando el término lineal extraído vía el residuo:
    $$ \prod_{n=1}^{\infty} \lambda_n := \exp\left( -\text{Res}_{s=0} \frac{\zeta_\Lambda(s)}{s^2} \right) $$

    \item \textbf{Suma Zeta-regularizada:} De manera análoga, la suma regularizada de los elementos de la secuencia (o de sus potencias) se define evaluando la función zeta en enteros negativos. La suma de los propios elementos de la secuencia corresponde a evaluar la función en $s = -1$:
    $$ \sum_{n=1}^{\infty} \lambda_n := \zeta_\Lambda(-1) $$
    Más generalmente, la suma regularizada de las $k$-ésimas potencias de los elementos de la secuencia se define como:
    $$ \sum_{n=1}^{\infty} \lambda_n^k := \zeta_\Lambda(-k) $$
    Un ejemplo célebre es la suma de todos los números naturales, donde $\Lambda = \mathbb{N}$ y la función es la Zeta de Riemann $\zeta(s)$. La suma regularizada es $\sum_{n=1}^\infty n = \zeta(-1) = -1/12$.
\end{itemize}

La fórmula maestra $\prod \lambda_i = \exp(-\zeta_\Lambda'(0))$ no es, por tanto, un mero truco algebraico, sino la consecuencia directa de identificar el determinante funcional (o producto infinito) con la derivada en el origen de la traza regularizada del operador (o suma de Dirichlet). Esta definición preserva la multiplicatividad en la mayoría de los casos, aunque, como veremos en clases posteriores, puede dar lugar a profundas \textit{anomalías multiplicativas} cuando se regularizan productos de productos.

\section{Propiedades algebraicas básicas}

Una vez establecido el formalismo analítico para asignar valores finitos a sumas y productos divergentes, es imperativo investigar cómo estas operaciones regularizadas interactúan con las estructuras algebraicas clásicas. A primera vista, podría pensarse que la introducción de la continuación analítica —una herramienta global y altamente no lineal del análisis complejo— destruiría las propiedades algebraicas elementales que poseen las sumas y los productos finitos. Sin embargo, la Zeta-Regularización demuestra una notable resiliencia estructural: hereda y generaliza varias propiedades fundamentales, aunque a menudo con matices analíticos profundos que carecen de análogo en el análisis real clásico.

A continuación, desarrollamos las tres propiedades algebraicas básicas: la linealidad de la suma regularizada, la multiplicatividad bajo particiones de conjuntos, y el fascinante comportamiento ante escalas (homogeneidad generalizada) del producto regularizado.

\subsection{Linealidad de la suma regularizada y aditividad espectral}

La propiedad de linealidad en el contexto de la Zeta-Regularización se manifiesta de dos formas distintas: la homogeneidad de la suma regularizada respecto a la multiplicación por escalares, y la aditividad estricta sobre uniones disjuntas.

Dada una secuencia $\Lambda = \{\lambda_n\}$ y una constante compleja $c \neq 0$, definimos la secuencia escalada $c\Lambda = \{c\lambda_n\}$. La función zeta asociada a esta nueva secuencia es:

$$ \zeta_{c\Lambda}(s) = \sum_{n=1}^{\infty} (c\lambda_n)^{-s} = c^{-s} \sum_{n=1}^{\infty} \lambda_n^{-s} = c^{-s} \zeta_\Lambda(s) $$

Si evaluamos esta relación para la suma regularizada, la cual se define evaluando la función zeta en $s = -1$, obtenemos:

$$ \sum_{n=1}^{\infty} (c\lambda_n) := \zeta_{c\Lambda}(-1) = c^{-(-1)} \zeta_\Lambda(-1) = c \sum_{n=1}^{\infty} \lambda_n $$

Esto demuestra que la suma Zeta-regularizada es estrictamente lineal (homogénea de grado 1) respecto a la multiplicación por escalares, comportándose exactamente como una suma finita clásica.

Por otro lado, la aditividad sobre particiones es una consecuencia directa de la definición de la serie de Dirichlet. Si el conjunto índice se puede particionar en dos subconjuntos disjuntos $A$ y $B$ tales que $\Lambda = A \sqcup B$, la función zeta total es simplemente la suma de las funciones zeta de las partes:

$$ \zeta_\Lambda(s) = \zeta_A(s) + \zeta_B(s) $$

Esta linealidad espectral se preserva para cualquier evaluación, garantizando que la suma regularizada total es la suma de las regularizaciones individuales:

$$ \sum_{\lambda \in \Lambda} \lambda = \sum_{\alpha \in A} \alpha + \sum_{\beta \in B} \beta $$

\subsection{Particiones de conjuntos y multiplicatividad del producto}

Si la linealidad de la suma es directa, la propiedad correspondiente para el producto —la multiplicatividad bajo particiones de conjuntos— es igual de robusta y resulta ser una de las herramientas más poderosas en aplicaciones físicas, como la separación de modos en el efecto Casimir.

Supongamos nuevamente que $\Lambda = A \sqcup B$. La relación aditiva de las funciones zeta, $\zeta_\Lambda(s) = \zeta_A(s) + \zeta_B(s)$, se mantiene en todo el dominio de convergencia y, por el teorema de identidad analítica, en todo el plano complejo tras la continuación. 

Al derivar ambos lados respecto a $s$ y evaluar en $s=0$, obtenemos una relación puramente aditiva para las derivadas:

$$ \zeta_\Lambda'(0) = \zeta_A'(0) + \zeta_B'(0) $$

Dado que el producto regularizado se define mediante la función exponencial de estas derivadas, la aditividad en el argumento de la exponencial se transforma en multiplicatividad en el resultado final:

$$ \prod_{\lambda \in \Lambda} \lambda := \exp(-\zeta_\Lambda'(0)) = \exp(-\zeta_A'(0) - \zeta_B'(0)) = \exp(-\zeta_A'(0)) \exp(-\zeta_B'(0)) $$

Por lo tanto, se cumple estrictamente que:

$$ \prod_{\lambda \in \Lambda} \lambda = \left( \prod_{\alpha \in A} \alpha \right) \left( \prod_{\beta \in B} \beta \right) $$

\begin{itemize}
    \item \textbf{Implicación analítica:} Esta propiedad nos permite tomar un producto infinitamente divergente, fragmentarlo en sub-productos que también son individualmente divergentes, regularizar cada uno por separado mediante sus propias funciones zeta, y luego multiplicar los resultados finitos. El orden en que se agrupen los términos (siempre que la partición sea bien definida y las funciones zeta parciales admitan continuación) no altera el resultado final.
\end{itemize}

\subsection{Comportamiento ante escalas (Homogeneidad generalizada)}

La propiedad más contraintuitiva y analíticamente rica de la Zeta-Regularización surge al investigar el comportamiento del producto infinito bajo la multiplicación de todos sus términos por una constante $a > 0$. En el análisis clásico, si tenemos un conjunto finito de $N$ elementos, la extracción de la constante $a$ del producto sigue la regla de homogeneidad: $\prod_{i=1}^N (a\lambda_i) = a^N \prod_{i=1}^N \lambda_i$. 

Veamos cómo se generaliza esto en el régimen regularizado. Partimos de la relación establecida previamente para la función zeta escalada:

$$ \zeta_{a\Lambda}(s) = a^{-s} \zeta_\Lambda(s) $$

Para encontrar el producto regularizado de la secuencia escalada, necesitamos la derivada de esta función evaluada en $s=0$. Aplicamos la regla del producto para la derivación, recordando que $a^{-s} = e^{-s \log a}$:

$$ \zeta_{a\Lambda}'(s) = \frac{d}{ds} \left( e^{-s \log a} \zeta_\Lambda(s) \right) = -\log(a) e^{-s \log a} \zeta_\Lambda(s) + e^{-s \log a} \zeta_\Lambda'(s) $$

Evaluando esta expresión en $s=0$, y dado que $e^0 = 1$, obtenemos:

$$ \zeta_{a\Lambda}'(0) = -\log(a) \zeta_\Lambda(0) + \zeta_\Lambda'(0) $$

Ahora, aplicamos la fórmula maestra para el producto regularizado:

$$ \prod_{n=1}^{\infty} (a\lambda_n) = \exp(-\zeta_{a\Lambda}'(0)) = \exp\left( \log(a) \zeta_\Lambda(0) - \zeta_\Lambda'(0) \right) $$

Utilizando las propiedades de la exponencial, esto se factoriza como:

$$ \prod_{n=1}^{\infty} (a\lambda_n) = \exp(\log(a) \zeta_\Lambda(0)) \exp(-\zeta_\Lambda'(0)) = a^{\zeta_\Lambda(0)} \prod_{n=1}^{\infty} \lambda_n $$

\begin{enumerate}[a)]
    \item \textbf{La generalización del "número de elementos":} La fórmula $a^{\zeta_\Lambda(0)}$ es el análogo exacto de $a^N$ en el caso finito. Esto revela un hecho analítico profundo: el valor de la función zeta evaluada en el origen, $\zeta_\Lambda(0)$, actúa como el \textit{número regularizado de elementos} en la secuencia $\Lambda$. 
    \item \textbf{Ejemplo paradigmático:} Si $\Lambda = \mathbb{N}$, la función zeta es la función Zeta de Riemann $\zeta(s)$. Sabemos que $\zeta(0) = -1/2$. Si escalamos todos los números naturales por una constante $a$, el producto regularizado no escala como $a^\infty$ (como sugeriría la divergencia clásica), ni como $a^0$, sino como $a^{-1/2}$. Es decir, la secuencia de los números naturales se comporta algebraicamente como si tuviera "$-1/2$ elementos".
    \item \textbf{Dependencia de la escala en el determinante funcional:} En geometría espectral, si escalamos la métrica de una variedad de Riemann por un factor $a^2$, los valores propios del Laplaciano se escalan por $a^{-2}$. La fórmula de homogeneidad generalizada dicta exactamente cómo cambia el determinante funcional del operador bajo cambios de escala de la geometría subyacente, siendo $\zeta_\Delta(0)$ un invariante espectral crucial que codifica la dimensión y la topología del espacio.
\end{enumerate}

En resumen, mientras que la suma regularizada preserva la linealidad clásica de manera trivial, el producto regularizado reemplaza el conteo entero clásico de elementos por la evaluación analítica $\zeta_\Lambda(0)$, demostrando que la Zeta-Regularización no destruye el álgebra subyacente, sino que la eleva a un marco donde la aritmética finita es gobernada por la topología y el análisis complejo.

\section{Ejemplos clásicos y la Fórmula de Lerch}

Habiendo establecido el marco analítico general de la Zeta-Regularización —la construcción de la función zeta espectral, su continuación analítica, la fórmula maestra y sus propiedades algebraicas—, es momento de poner en marcha la maquinaria sobre el ejemplo más fundamental y paradigmático: la secuencia de los números naturales $\Lambda = \mathbb{N} = \{1, 2, 3, \dots\}$. Este caso, históricamente el primero en ser abordado por Euler, no solo ilustra con claridad cristalina cada paso del procedimiento, sino que revela una conexión profunda y sorprendente con la función Gamma de Euler y el teorema de factorización de Weierstrass. El resultado final, conocido como la \textit{Fórmula de Lerch}, es una de las joyas del análisis clásico y el punto de partida de innumerables generalizaciones.

\subsection{El producto de todos los números naturales: $\sqrt{2\pi}$}

Consideremos la secuencia $\Lambda = \mathbb{N}$, cuya función zeta asociada es la célebre función Zeta de Riemann:

$$ \zeta(s) = \sum_{n=1}^{\infty} n^{-s}, \qquad \Re(s) > 1 $$

La abscisa de convergencia es $\sigma_c = 1$. Nuestro objetivo es calcular el producto regularizado de todos los números naturales:

$$ P = \prod_{n=1}^{\infty} n = \exp(-\zeta'(0)) $$

Para ello, necesitamos evaluar la derivada de $\zeta(s)$ en $s=0$, un punto que se encuentra estrictamente a la izquierda de la abscisa de convergencia. La herramienta clave es la \textbf{ecuación funcional} de la función Zeta de Riemann, válida por continuación analítica para todo $s \in \mathbb{C}$:

$$ \zeta(s) = 2^s \pi^{s-1} \sin\left(\frac{\pi s}{2}\right) \Gamma(1-s) \zeta(1-s) $$

Para extraer $\zeta'(0)$, tomamos el logaritmo de ambos lados (trabajando formalmente en un entorno de $s=0$ donde todas las funciones involucradas son no nulas o tienen comportamiento controlado) y derivamos. Sin embargo, es más instructivo proceder mediante la expansión asintótica en serie de Taylor alrededor de $s=0$. Analicemos cada factor de la ecuación funcional:

\begin{enumerate}[a)]
    \item \textbf{Factor exponencial:} $2^s \pi^{s-1} = \frac{1}{\pi} e^{s \log(2\pi)} = \frac{1}{\pi}\left(1 + s \log(2\pi) + \mathcal{O}(s^2)\right)$.
    
    \item \textbf{Factor trigonométrico:} $\sin\left(\frac{\pi s}{2}\right) = \frac{\pi s}{2} - \frac{(\pi s)^3}{48} + \dots = \frac{\pi s}{2}\left(1 + \mathcal{O}(s^2)\right)$.
    
    \item \textbf{Factor Gamma:} $\Gamma(1-s) = 1 + \gamma s + \mathcal{O}(s^2)$, donde $\gamma$ es la constante de Euler-Mascheroni.
    
    \item \textbf{Factor zeta reflejado:} $\zeta(1-s)$. Cerca de $s=0$, el argumento $1-s$ se acerca a $1$, donde $\zeta$ tiene un polo simple con residuo $1$. La expansión de Laurent es $\zeta(1-s) = \frac{1}{-s} + \gamma + \mathcal{O}(s) = -\frac{1}{s} + \gamma + \mathcal{O}(s)$.
\end{enumerate}

Multiplicando todos los factores y recogiendo términos, la ecuación funcional nos da:

$$ \zeta(s) = \frac{1}{\pi}\left(1 + s\log(2\pi) + \dots\right) \cdot \frac{\pi s}{2}\left(1 + \dots\right) \cdot \left(1 + \gamma s + \dots\right) \cdot \left(-\frac{1}{s} + \gamma + \dots\right) $$

Simplificando el producto de los términos principales:

$$ \zeta(s) = \frac{s}{2} \cdot \left(-\frac{1}{s}\right) + \mathcal{O}(s) = -\frac{1}{2} + \mathcal{O}(s) $$

Esto confirma el valor bien conocido $\zeta(0) = -1/2$, que actúa como el ``número regularizado'' de elementos en $\mathbb{N}$. Para obtener la derivada, necesitamos el término lineal en $s$. Tras una multiplicación cuidadosa de las expansiones y recolección del coeficiente de $s$, se obtiene:

$$ \zeta(s) = -\frac{1}{2} - \frac{1}{2}\log(2\pi) \cdot s + \mathcal{O}(s^2) $$

De donde leemos directamente:

$$ \zeta'(0) = -\frac{1}{2}\log(2\pi) $$

Finalmente, aplicando la fórmula maestra:

$$ P = \prod_{n=1}^{\infty} n = \exp(-\zeta'(0)) = \exp\left(\frac{1}{2}\log(2\pi)\right) = \sqrt{2\pi} $$

Este resultado, originalmente obtenido por Euler mediante métodos formales mucho antes de que se inventara la continuación analítica, queda ahora rigurosamente justificado. El producto infinito de todos los números naturales, que clásicamente diverge hacia $+\infty$, posee el valor regularizado finito $\sqrt{2\pi}$. Es notable que este mismo valor aparece en la fórmula de Stirling, en la normalización de la distribución gaussiana y en el volumen de la esfera unidad, revelando una unidad estructural profunda en las matemáticas.

\subsection{La función Zeta de Hurwitz}

Para generalizar el resultado anterior y llegar a la Fórmula de Lerch, necesitamos introducir un objeto analítico más rico: la \textbf{función Zeta de Hurwitz}. Dado un parámetro real $x > 0$, se define como:

$$ \zeta(s, x) = \sum_{n=0}^{\infty} \frac{1}{(n+x)^s}, \qquad \Re(s) > 1 $$

Esta función generaliza la Zeta de Riemann, ya que $\zeta(s, 1) = \zeta(s)$. La secuencia asociada es $\Lambda_x = \{x, x+1, x+2, \dots\}$, que representa un desplazamiento aditivo de los enteros no negativos. La abscisa de convergencia sigue siendo $\sigma_c = 1$ para todo $x > 0$.

La función $\zeta(s, x)$ admite una continuación meromorfa a todo el plano complejo con un único polo simple en $s=1$ de residuo $1$. Lo crucial para la regularización es que $s=0$ es un punto regular, y los valores de la función y su derivada en el origen admiten expresiones cerradas en términos de la función Gamma:

\begin{itemize}
    \item \textbf{Valor en el origen:}
    $$ \zeta(0, x) = \frac{1}{2} - x $$
    
    \item \textbf{Derivada en el origen:}
    $$ \zeta'(0, x) = \log\left(\frac{\Gamma(x)}{\sqrt{2\pi}}\right) $$
\end{itemize}

La primera identidad generaliza el resultado $\zeta(0) = \zeta(0, 1) = 1/2 - 1 = -1/2$. La segunda identidad es la piedra angular de la Fórmula de Lerch y merece una derivación detallada.

\subsubsection{Derivación de $\zeta'(0, x)$}

La demostración se apoya en la representación integral de la función Zeta de Hurwitz mediante la transformada de Mellin:

$$ \zeta(s, x) = \frac{1}{\Gamma(s)} \int_0^{\infty} \frac{t^{s-1} e^{-xt}}{1 - e^{-t}} \, dt, \qquad \Re(s) > 1, \quad x > 0 $$

Para acceder a la región de $s=0$, se separa la integral en dos partes: una cercana al origen ($t \in [0, 1]$) donde se expande el integrando en serie, y una cola ($t \in [1, \infty)$) que ya es convergente para todo $s$. La expansión del núcleo $\frac{e^{-xt}}{1-e^{-t}}$ alrededor de $t=0$ involucra los polinomios de Bernoulli, y tras una regularización cuidadosa del polo en $s=1$ y la evaluación del término finito, se llega a la expresión:

$$ \zeta'(0, x) = \log\Gamma(x) - \frac{1}{2}\log(2\pi) = \log\left(\frac{\Gamma(x)}{\sqrt{2\pi}}\right) $$

Esta identidad revela que la derivada en el origen de la función Zeta de Hurwitz codifica, esencialmente, la función Gamma de Euler.

\subsection{La Fórmula de Lerch}

Con los ingredientes anteriores, la \textbf{Fórmula de Lerch} emerge como una consecuencia inmediata y elegante de la fórmula maestra. El producto Zeta-regularizado de la secuencia desplazada $\Lambda_x = \{x, x+1, x+2, \dots\}$ se define como:

$$ \prod_{n=0}^{\infty} (n+x) := \exp(-\zeta'(0, x)) $$

Sustituyendo la expresión obtenida para $\zeta'(0, x)$:

$$ \prod_{n=0}^{\infty} (n+x) = \exp\left(-\log\left(\frac{\Gamma(x)}{\sqrt{2\pi}}\right)\right) = \frac{\sqrt{2\pi}}{\Gamma(x)} $$

Esta es la \textbf{Fórmula de Lerch}, válida para todo $x > 0$. Escribiéndola de manera prominente:

$$ \boxed{\prod_{n=0}^{\infty} (n+x) = \frac{\sqrt{2\pi}}{\Gamma(x)}} $$

\begin{itemize}
    \item \textbf{Caso particular $x=1$:} Se recupera el producto de todos los números naturales:
    $$ \prod_{n=1}^{\infty} n = \prod_{n=0}^{\infty} (n+1) = \frac{\sqrt{2\pi}}{\Gamma(1)} = \frac{\sqrt{2\pi}}{1} = \sqrt{2\pi} $$
    
    \item \textbf{Caso particular $x=1/2$:} Se obtiene el producto de los semienteros positivos:
    $$ \prod_{n=0}^{\infty} \left(n + \frac{1}{2}\right) = \frac{\sqrt{2\pi}}{\Gamma(1/2)} = \frac{\sqrt{2\pi}}{\sqrt{\pi}} = \sqrt{2} $$
\end{itemize}

\subsection{Conexión con el producto infinito de Weierstrass}

La Fórmula de Lerch no existe en el vacío: es la versión ``renormalizada'' del célebre \textbf{teorema de factorización de Weierstrass} para la función Gamma. El producto de Weierstrass establece que:

$$ \frac{1}{\Gamma(x)} = x \, e^{\gamma x} \prod_{n=1}^{\infty} \left(1 + \frac{x}{n}\right) e^{-x/n} $$

donde $\gamma = \lim_{N \to \infty} \left(\sum_{n=1}^{N} \frac{1}{n} - \log N\right) \approx 0.5772$ es la constante de Euler-Mascheroni. 

Para apreciar la conexión, reescribamos el producto de Weierstrass agrupando los factores de manera diferente. Multiplicando ambos lados por $\sqrt{2\pi}$:

$$ \frac{\sqrt{2\pi}}{\Gamma(x)} = \sqrt{2\pi} \cdot x \, e^{\gamma x} \prod_{n=1}^{\infty} \left(1 + \frac{x}{n}\right) e^{-x/n} $$

El lado izquierdo es precisamente la Fórmula de Lerch. El lado derecho es el producto de Weierstrass ``vestido'' con factores de convergencia $e^{-x/n}$. La relación profunda entre ambos es la siguiente:

\begin{enumerate}[a)]
    \item \textbf{El producto de Weierstrass es convergente:} Los factores $e^{-x/n}$ son \textit{factores de convergencia} que se introducen artificialmente para forzar la convergencia del producto infinito. Sin ellos, el producto $\prod (1 + x/n)$ diverge.
    
    \item \textbf{La Fórmula de Lerch es divergente pero regularizada:} El producto $\prod_{n=0}^{\infty} (n+x)$ diverge clásicamente, pero la Zeta-Regularización le asigna el valor finito $\sqrt{2\pi}/\Gamma(x)$. 
    
    \item \textbf{Equivalencia analítica:} Ambas fórmulas expresan la misma información matemática, pero en registros distintos. El producto de Weierstrass es una identidad en el análisis clásico (convergencia ordinaria). La Fórmula de Lerch es una identidad en el análisis regularizado (convergencia zeta). Los factores de convergencia $e^{-x/n}$ del producto de Weierstrass son precisamente los ``contra-términos'' que la regularización zeta sustrae automáticamente a través de la continuación analítica.
\end{enumerate}

Esta equivalencia puede verse como un principio general: \textit{la Zeta-Regularización produce resultados que son consistentes con, y a menudo equivalentes a, las factorizaciones de Weierstrass, pero sin necesidad de introducir factores de convergencia ad hoc}. La continuación analítica hace el trabajo de regularización de manera intrínseca y canónica.

\subsection{Ejemplos derivados y la Fórmula de Lerch generalizada}

La Fórmula de Lerch y sus propiedades algebraicas generan una familia rica de identidades reguladas. A continuación se presentan algunos ejemplos notables.

\subsubsection{Productos de progresiones aritméticas}

Utilizando la propiedad de homogeneidad generalizada ($\prod a\lambda_n = a^{\zeta_\Lambda(0)} \prod \lambda_n$) y la partición de conjuntos, podemos calcular productos sobre subconjuntos de $\mathbb{N}$.

\begin{itemize}
    \item \textbf{Producto de los números pares:} La secuencia $\{2, 4, 6, \dots\}$ es $2\mathbb{N}$. Por la propiedad de escala, con $\zeta(0) = -1/2$:
    $$ \prod_{n=1}^{\infty} (2n) = 2^{\zeta(0)} \prod_{n=1}^{\infty} n = 2^{-1/2} \cdot \sqrt{2\pi} = \sqrt{\pi} $$
    
    \item \textbf{Producto de los números impares:} Como $\mathbb{N} = \{2n\} \sqcup \{2n-1\}$, la multiplicatividad por particiones da:
    $$ \prod_{n=1}^{\infty} (2n-1) = \frac{\prod_{n=1}^{\infty} n}{\prod_{n=1}^{\infty} (2n)} = \frac{\sqrt{2\pi}}{\sqrt{\pi}} = \sqrt{2} $$
\end{itemize}

\subsubsection{La Fórmula de Lerch generalizada para formas cuadráticas}

Una generalización notable, debida originalmente al propio Lerch y posteriormente extendida por Kurokawa y Wakayama, involucra productos sobre expresiones cuadráticas. Para $x > 0$ e $y \in \mathbb{R}$, se tiene:

$$ \prod_{n=0}^{\infty} \left((n+x)^2 + y^2\right) = \frac{2\pi}{\Gamma(x+iy)\Gamma(x-iy)} $$

Esta identidad se deriva de la observación de que la función zeta asociada a la secuencia $\{(n+x)^2 + y^2\}$ puede descomponerse en una combinación de funciones Zeta de Hurwitz evaluadas en $x+iy$ y $x-iy$. Nótese que para $y=0$ se recupera la Fórmula de Lerch clásica (elevada al cuadrado).

\subsubsection{El producto $\prod n^n$ y la constante de Glaisher-Kinkelin}

Un ejemplo que requiere ir más allá de $\zeta'(0)$ es el producto $\prod_{n=1}^{\infty} n^n$. La función zeta asociada es $\sum_{n=1}^{\infty} (n^n)^{-s} = \sum_{n=1}^{\infty} n^{-ns}$, pero el enfoque más natural es notar que:

$$ \log \prod_{n=1}^{\infty} n^n = \sum_{n=1}^{\infty} n \log n = -\left.\frac{d}{ds}\sum_{n=1}^{\infty} n^{-s}\right|_{s=-1} \cdot (-1) $$

En realidad, se utiliza directamente que $\sum n \log n$ se regulariza como $-\zeta'(-1)$, de modo que:

$$ \prod_{n=1}^{\infty} n^n = \exp(-\zeta'(-1)) = A \, e^{-1/12} $$

donde $A \approx 1.2824271\dots$ es la \textbf{constante de Glaisher-Kinkelin}. Este ejemplo muestra que la maquinaria zeta-regularizada no se limita a $s=0$; evaluaciones en otros enteros negativos producen constantes matemáticas de profundo interés.

\section{Series de potencias finitas y la función Zeta de Hurwitz}

Hasta ahora, hemos utilizado la función Zeta de Riemann para regularizar sumas sobre los números naturales. Sin embargo, en aplicaciones tanto matemáticas como físicas, es frecuente encontrarse con secuencias desplazadas, progresiones aritméticas que no inician en cero, o sumas finitas truncadas por un corte (cut-off) energético. Para abordar estas situaciones con el mismo rigor analítico, debemos introducir una generalización natural: la \textbf{función Zeta de Hurwitz}. Esta función no solo nos permite evaluar sumas infinitas desplazadas, sino que proporciona la herramienta algebraica exacta para descomponer sumas finitas y aislar sus partes divergentes mediante expansiones asintóticas.

\subsection{Definición y continuación analítica de la función Zeta de Hurwitz}

Dado un número complejo $a$ con $\Re(a) > 0$, la función Zeta de Hurwitz se define inicialmente mediante la serie de Dirichlet:

$$ \zeta(s, a) = \sum_{n=0}^{\infty} \frac{1}{(n+a)^s} $$

Esta serie converge absolutamente para $\Re(s) > 1$. Al igual que la función Zeta de Riemann, $\zeta(s, a)$ admite una continuación analítica a todo el plano complejo $s \in \mathbb{C}$, excepto por un polo simple en $s=1$ con residuo $1$. 

La conexión con la función de Riemann es inmediata: $\zeta(s, 1) = \zeta(s)$. No obstante, la dependencia en el parámetro $a$ enriquece enormemente la estructura analítica. En particular, los valores de $\zeta(s, a)$ en los enteros negativos no triviales están íntimamente ligados a los \textbf{polinomios de Bernoulli} $B_k(x)$. Para cualquier entero $m \ge 0$, se cumple la identidad exacta:

$$ \zeta(-m, a) = -\frac{B_{m+1}(a)}{m+1} $$

Esta fórmula es la piedra angular para evaluar sumas de potencias de secuencias desplazadas. Por ejemplo, para $m=0$ y $m=1$, obtenemos respectivamente:

\begin{itemize}
    \item $\zeta(0, a) = \frac{1}{2} - a$
    \item $\zeta(-1, a) = -\frac{1}{2}\left(a^2 - a + \frac{1}{6}\right)$
\end{itemize}

\subsection{Evaluación de sumas finitas y desplazadas}

La verdadera potencia de la función Zeta de Hurwitz en el contexto de la regularización radica en su capacidad para expresar \textit{sumas finitas} de potencias como la diferencia de dos funciones Zeta de Hurwitz evaluadas en enteros negativos. 

Supongamos que deseamos evaluar la suma finita de las $m$-ésimas potencias de una secuencia desplazada, truncada en un índice $N$:

$$ S_N(m, a) = \sum_{n=0}^{N-1} (n+a)^m $$

Podemos reescribir esta suma finita como la diferencia entre una suma infinita (desde $n=0$) y una suma infinita desplazada (desde $n=N$):

$$ S_N(m, a) = \sum_{n=0}^{\infty} (n+a)^m - \sum_{n=N}^{\infty} (n+a)^m $$

En el análisis clásico, ambas sumas divergen si $m \ge 0$. Sin embargo, en el esquema de Zeta-Regularización, aplicamos la continuación analítica a cada término por separado. El primer término es simplemente $\zeta(-m, a)$. Para el segundo término, realizamos un cambio de índice $k = n - N$, de modo que la suma se convierte en $\sum_{k=0}^{\infty} (k + N + a)^m$, la cual es exactamente $\zeta(-m, N+a)$. 

Por lo tanto, obtenemos la identidad fundamental para sumas finitas regularizadas:

$$ \sum_{n=0}^{N-1} (n+a)^m = \zeta(-m, a) - \zeta(-m, N+a) $$

Esta ecuación es exacta bajo la regularización zeta. Nos dice que la suma finita truncada en $N$ no es más que el valor regularizado infinito $\zeta(-m, a)$ menos la "cola" infinita que comienza en $N$, la cual está codificada en $\zeta(-m, N+a)$.

\subsection{Expansión asintótica para grandes valores del corte}

Para entender cómo la regularización extrae valores finitos en presencia de divergencias (un mecanismo que estudiaremos a fondo en aplicaciones físicas), debemos analizar el comportamiento de la "cola" $\zeta(-m, N+a)$ cuando el corte $N$ es muy grande ($N \to \infty$). 

Para ello, utilizamos la expansión asintótica de la función Zeta de Hurwitz para $a \to \infty$ con $s$ fijo, la cual se deriva directamente de la fórmula de sumación de Euler-Maclaurin:

$$ \zeta(s, a) \sim \frac{a^{1-s}}{s-1} + \frac{1}{2}a^{-s} - \sum_{k=1}^{\infty} \frac{B_{2k}}{(2k)!} s(s+1)\cdots(s+2k-2) a^{-s-2k+1} $$

Si evaluamos esta expansión en un entero negativo $s = -m$ (donde $m \ge 0$), la serie asintótica se convierte en un polinomio exacto en $a$ (coincidiendo con el polinomio de Bernoulli $-B_{m+1}(a)/(m+1)$). Despejando los términos, la expansión para la "cola" truncada toma la forma:

$$ \zeta(-m, N+a) \sim -\frac{(N+a)^{m+1}}{m+1} - \frac{1}{2}(N+a)^m + \sum_{k=1}^{\lfloor m/2 \rfloor} \frac{B_{2k}}{(2k)!} \frac{(m+1)!}{(m+2-2k)!} (N+a)^{m+1-2k} $$

\subsubsection{El mecanismo de cancelación de divergencias}

La expansión anterior revela la estructura analítica de las sumas truncadas. Si sustituimos esta expansión en la identidad de la suma finita, obtenemos:

$$ \sum_{n=0}^{N-1} (n+a)^m = \zeta(-m, a) + \frac{(N+a)^{m+1}}{m+1} + \frac{1}{2}(N+a)^m - \dots $$

Aquí yace el corazón del método de regularización por sustracción de estados de referencia:

\begin{enumerate}[a)]
    \item \textbf{Términos Divergentes (Polinomios en $N$):} Los términos $\frac{N^{m+1}}{m+1}$, $\frac{N^m}{2}$, etc., crecen sin límite cuando $N \to \infty$. En el contexto físico, estos términos representan la energía del vacío no observable, la cual depende del volumen del sistema o del corte ultravioleta, pero no de la geometría específica ni de las condiciones de frontera.
    \item \textbf{Término Regularizado (Constante independiente de $N$):} El término $\zeta(-m, a)$ es completamente independiente del corte $N$. Este es el valor finito, intrínseco y observable que la Zeta-Regularización asigna a la suma.
    \item \textbf{Cancelación Exacta:} En problemas físicos (como el efecto Casimir o la magnetización de materiales cuánticos), se introduce un estado de referencia (por ejemplo, el vacío sin placas o el gas de electrones sin campo magnético) cuya energía de referencia se calcula mediante una integral continua. Matemáticamente, esta integral continua genera exactamente los mismos términos polinómicos en $N$ que la expansión asintótica de $\zeta(-m, N+a)$. Al restar la energía de referencia a la suma discreta, todos los términos divergentes (las potencias positivas de $N$) se cancelan analítica y exactamente, dejando únicamente el valor regularizado $\zeta(-m, a)$.
\end{enumerate}

\subsection{Ejemplo ilustrativo: Suma de enteros con corte}

Para consolidar esta idea, consideremos la suma de los primeros $N$ números naturales, es decir, $a=1$ y $m=1$:

$$ \sum_{n=1}^{N} n = \zeta(-1, 1) - \zeta(-1, N+1) $$

Sabemos que $\zeta(-1, 1) = \zeta(-1) = -1/12$. Usando la expansión asintótica para $\zeta(-1, N+1)$ con $m=1$:

$$ \zeta(-1, N+1) = -\frac{(N+1)^2}{2} - \frac{1}{2}(N+1) + \frac{B_2}{2}(1) = -\frac{N^2}{2} - N - \frac{1}{2} + \frac{1}{12} $$

Sustituyendo en la identidad de la suma finita:

$$ \sum_{n=1}^{N} n = -\frac{1}{12} - \left( -\frac{N^2}{2} - N - \frac{1}{2} + \frac{1}{12} \right) = \frac{N^2}{2} + \frac{N}{2} = \frac{N(N+1)}{2} $$

Hemos recuperado la fórmula clásica de la suma de una progresión aritmética. Sin embargo, la perspectiva analítica es profunda: la fórmula clásica es el resultado de tomar la suma infinita regularizada ($-1/12$) y sumarle la "cola" divergente expandida asintóticamente. 

Cuando en física nos enfrentamos a una suma que no tiene una forma cerrada simple, o cuando el corte $N$ depende de parámetros físicos continuos (como la intensidad de un campo magnético o la distancia entre placas), no podemos usar la fórmula cerrada $\frac{N(N+1)}{2}$. En su lugar, \textit{debemos} usar la expansión asintótica de $\zeta(s, N+a)$ para aislar los términos divergentes, cancelarlos contra la integral de referencia, y extraer el término constante $\zeta(s, a)$, el cual constituye la predicción física regularizada.
